% !TeX root = closed-systems-from-comparison-completeness.tex
% ============================================================
\chapter{Conditional Completion Ledger}
\label{app:conditional-completion}
% ============================================================

\section{Extracted structural content}
\label{sec:cond-reduced}

This ledger records entries that are structurally reduced but not yet
numerically fixed by closed transport alone.
In the Chapter~20/21 terminology, such entries are called
\emph{extracted}.
No additional theorem-level claim is introduced here; this chapter is
status bookkeeping for downstream numerical closure. The state-side
reconstruction extension of \cref{ch:dynamics-program} has its own
status ledger and open-problem list; it is not part of the
Chapter~14--22 numerical-closure audit recorded here.

Within the theorem stack recorded here, the extracted entry is:
\begin{enumerate}[label=\textup{(\roman*)}]
\item
  full representation-content classification reduced to irreducible unitary
  representation data of
  \(U(1)\times SU(2)\times SU(3)\)
  (\cref{thm:matter-spectrum-reduction,cor:matter-unique-forced-output}).
\end{enumerate}

\section{Open quantitative boundary data}
\label{sec:cond-open}

This section records quantitatively open entries that are not determined by
the theorem-level closure chain of Chapters~1--22.
Unless further closure theorems are proved, these entries remain explicit
boundary data.

After the global closure audit of Chapters~14--22, the ledger is typed in
three layers.

\textbf{Layer A (internal bridge data discharged).}
The scalar bilinear source is forced by polarization
(\cref{thm:scalar-pairing}), and the closed-system quartic mass coordinate is
fixed by the intrinsic readout bridge
(\cref{thm:numbers-mass-calibration}).

\textbf{Layer B (properties and data discharged by theorem).}
The explicit depth-anchor package \cref{def:numbers-depth-anchors} is
optional: it is discharged whenever the ratio/cubic calibration criterion
of
\cref{thm:numbers-anchor-closure-61835,cor:numbers-anchor-closure-discharges}
is available.

The shell-completeness property \cref{def:shell-completeness} is optional:
it is discharged whenever the signed-permutation quadratic criterion
\cref{thm:shell-completeness-from-symmetry,cor:shell-completeness-discharged}
is available, in particular via the Chapter~19 transport-derived route
\cref{thm:shell-symmetry-from-transport,cor:shell-completeness-from-transport},
with finite in-model verification by
\cref{thm:shell-interface-checklist}.

The one-reference shell closure data
\cref{def:shell-defect-intrinsic,def:shell-anchor} are no longer mandatory:
they are discharged whenever the two-reference algebraic closure criterion
of \cref{thm:shell-two-reference-closure,cor:shell-two-reference-discharges}
is available.

\textbf{Layer C (residual open quantitative-status entries).}
\begin{enumerate}[label=\textup{(\roman*)}]
\item
  explicit real-level form of the scalar quadratic channel
  \(Q_{\mathbb R}\)
  (\cref{rem:two-levels-Q,rem:spectral-scalar-id-scope});
\item
  shell-completeness status (explicit property declaration vs verified
  signed-permutation discharge criterion, including the Chapter~19
  transport-derived route) for data packages where
  \cref{thm:shell-completeness-from-symmetry,thm:shell-symmetry-from-transport,thm:shell-interface-checklist}
  is not established;
\item
  dimensionless shell-splitting constant status
  (internal determination vs irreducible \(\gamma\)) when neither the
  intrinsic-defect route nor two-reference algebraic closure is available;
\item
  absolute mass-scale normalization (open \(\alpha\)) when no calibration
  route fixes \(\alpha\);
\item
  generation-level realization details;
\item
  low-energy coupling-constant values.
\end{enumerate}

These entries remain outside the theorem-level closure recorded here and are tracked as
explicit open boundary data.

\section{Formal status ledger}
\label{sec:cond-status-ledger}

This section lists every global item that carries non-theorem-level load and
records its strongest status in this ledger.
The words ``datum'', ``property'', and ``principle'' are used deliberately:
a datum is retained only as boundary data, a property may be verified or
discharged theoremically, and the standing principle is the admissibility rule
that governs the whole construction.

\begin{enumerate}[label=\textup{(H\arabic*)}]
\item
  \textbf{Closed-world admissibility.}
  \Cref{sp:closed-world} is the sole standing principle.  It is not a hidden
  geometric, probabilistic, or dynamical postulate; it is the rule that only
  internally reconstructed comparison content is admissible.
\item
  \textbf{Scalar pairing bridge.}
  \Cref{thm:scalar-pairing} proves that the bilinear source needed to pass
  from the canonical scalar channel to a real quadratic form is forced by
  polarization; it is no longer retained as boundary data.
\item
  \textbf{Closed-system mass-coordinate bridge.}
  \Cref{thm:numbers-mass-calibration} proves that the intrinsic quartic
  coefficient-energy readout is the closed-system mass coordinate.  Only the
  external identification of that internal coordinate with a measured mass
  convention lies outside the internal theorem stack.
\item
  \textbf{Depth-anchor package.}
  \Cref{def:numbers-depth-anchors} is a package, not a condition.  It is
  stipulated directly only for data packages that choose explicit \((6,18,35)\)
  reference depths, and it is recovered by the ratio/cubic route
  \cref{thm:numbers-anchor-closure-61835,cor:numbers-anchor-closure-discharges}.
\item
  \textbf{Shell completeness.}
  \Cref{def:shell-completeness} is a property, not a condition.  It is
  proved by the signed-permutation quadratic criterion and, in transport-realized settings,
  by the transport-to-shell route
  \cref{thm:shell-completeness-from-symmetry,thm:shell-symmetry-from-transport,cor:shell-completeness-from-transport}.
\item
  \textbf{One-reference defect data.}
  \Cref{def:shell-defect-intrinsic} is one-reference shell bookkeeping
  data only.  It is unnecessary whenever the two-reference algebraic
  closure theorem applies.
\item
  \textbf{One-reference anchor data.}
  \Cref{def:shell-anchor} is likewise one-reference normalization data
  only.  It is discharged together with defect-evaluation data by
  \cref{thm:shell-two-reference-closure,cor:shell-two-reference-discharges}.
\end{enumerate}

\begin{definition}[Global closure-status vector]
\label{def:cond-status-vector}
The global closure-status vector used here is
\[
  \mathcal C_{\rm status}
  :=
  (P, B, M, A, S, R, E),
\]
where:
\begin{enumerate}[label=\textup{(\roman*)}]
\item \(P\) is the standing closed-world admissibility principle
      \cref{sp:closed-world};
\item \(B\) is the scalar-pairing bridge
      \cref{thm:scalar-pairing};
\item \(M\) is the closed-system mass-coordinate bridge
      \cref{thm:numbers-mass-calibration};
\item \(A\) is the depth-anchor package, either specified by
      \cref{def:numbers-depth-anchors} or recovered by
      \cref{thm:numbers-anchor-closure-61835};
\item \(S\) is shell completeness, either verified by
      \cref{thm:shell-completeness-from-symmetry,thm:shell-symmetry-from-transport}
      or retained as an open property;
\item \(R\) is one-reference shell bookkeeping, discharged when
      \cref{thm:shell-two-reference-closure} applies;
\item \(E\) is the residual observer-side semantic realization layer.
\end{enumerate}
\end{definition}

\begin{theorem}[Global closure audit normal form]
\label{thm:cond-closure-normal-form}
At the theorem strength recorded here, every global non-structural dependency in
Chapters~14--22 is represented by exactly one component of
\(\mathcal C_{\rm status}\).  Moreover:
\begin{enumerate}[label=\textup{(\roman*)}]
\item \(P\) is the only standing principle;
\item \(B\) and \(M\) are theorem-level bridges, not boundary data;
\item \(A\), \(S\), and \(R\) are either specified packages/properties or are
      discharged by the theorems named in \cref{def:cond-status-vector};
\item the only item not internalized by the mathematical closure stack is
      \(E\), the observer-side interpretation of closed-system coordinates as
      empirical readouts.
\end{enumerate}
\end{theorem}

\begin{proof}
The standing layer is unique by \cref{sp:closed-world} and by the
introductory scope statement, so all global admissibility load is accounted
for by \(P\).

For scalar structure, \cref{thm:scalar-pairing} constructs the bilinear source
by polarization from the real quadratic scalar channel.  Thus the
pairing bridge is represented by \(B\) and is theorem-level.  For mass
coordinates, \cref{thm:numbers-mass-calibration} identifies the internal
quartic coefficient-energy readout with the closed-system mass coordinate;
therefore the corresponding load is represented by \(M\) and is also
theorem-level.

For numerical anchors, \cref{def:numbers-depth-anchors} records the explicit
package, while \cref{thm:numbers-anchor-closure-61835,cor:numbers-anchor-closure-discharges}
recover and discharge it when the ratio/cubic route applies.  This is exactly
component \(A\).  For shell completeness,
\cref{def:shell-completeness} names the property and
\cref{thm:shell-completeness-from-symmetry,thm:shell-symmetry-from-transport,cor:shell-completeness-from-transport}
give the theorem-level verification routes; this is exactly component \(S\).
For one-reference shell bookkeeping,
\cref{def:shell-defect-intrinsic,def:shell-anchor} name the data and
\cref{thm:shell-two-reference-closure,cor:shell-two-reference-discharges}
discharge them under two-reference closure; this is exactly component \(R\).

The residual entries listed in \cref{sec:cond-open} either refine one of
\(A\), \(S\), or \(R\), or concern observer-side empirical interpretation.
\Cref{cor:mass-boundary,thm:matter-realization-factorization-criterion}
isolate that interpretation as a separate semantic realization problem,
which is component \(E\).  Hence every global non-structural dependency is
represented by one listed component, and the four claims follow.
\end{proof}

\begin{corollary}[No hidden global assumptions in the status normal form]
\label{cor:cond-no-hidden-global-assumptions}
After each global dependency is represented by its corresponding component
of \(\mathcal C_{\rm status}\), the manuscript has no unlisted global
quantitative assumption in Chapters~14--22.
\end{corollary}

\begin{proof}
By \cref{thm:cond-closure-normal-form}, every global non-structural dependency
belongs to exactly one component of \(\mathcal C_{\rm status}\).  Components
\(B\) and \(M\) are theorems; components \(A\), \(S\), and \(R\) are either
named explicitly or discharged by named theorems; component \(E\) is explicitly
semantic rather than mathematical.  Therefore no remaining global
quantitative assumption is hidden.
\end{proof}

\begin{remark}[Irreducible residue at recorded theorem strength]
In this status normal form,
\cref{def:numbers-depth-anchors,def:shell-completeness,def:shell-defect-intrinsic,def:shell-anchor}
carry no assumption-level status: each is either a defined package,
a theorem-proved property, or optional one-reference bookkeeping discharged by
a two-reference theorem.  The scalar-pairing and closed-system
mass-coordinate bridge theorems are proved internally in
\cref{thm:scalar-pairing,thm:numbers-mass-calibration}.
The only residue is semantic: an external observer may still ask whether the
closed-system mass coordinate is the measured mass used in a chosen empirical
protocol.  That identification is not an additional mathematical axiom of the
closed-system stack; it is an observer-side realization convention.
\end{remark}
